%% FPSAC 2027 extended abstract -- SKELETON v0 (2026-10-06)
%% Built by a writing sub-agent for Rick (AI agent). NOT FOR SUBMISSION. Robin submits.
%% Scope: memory/questions/q-fpsac-2027-writeup.md, STATUS v7.5 (Day 225 dream).
%% Grades: proofs/registry/hikita-star-dominance-support.json (registry is the truth, not prose).
%% Every theorem-like environment carries "% registry: <node-id>, grade: <grade>".
%% Template: official FPSAC2027.cls (downloaded 2026-10-06 from maths.universityofgalway.ie/fpsac2027/tex/).
%% Build: latexmk -pdf fpsac2027-draft.tex   (biblatex, backend=bibtex, loaded by the class)

\documentclass[submission]{FPSAC2027}

\newtheorem{theorem}{Theorem}[section]
\newtheorem{proposition}[theorem]{Proposition}
\newtheorem{lemma}[theorem]{Lemma}
\newtheorem{corollary}[theorem]{Corollary}
\newtheorem{conjecture}[theorem]{Conjecture}
\theoremstyle{definition}
\newtheorem{definition}[theorem]{Definition}
\newtheorem{example}[theorem]{Example}
\theoremstyle{remark}
\newtheorem{remark}[theorem]{Remark}

%% visible TODO marker (must stay visible in the PDF)
\newcommand{\todo}[1]{{\color{red}\textbf{[TODO:} #1\textbf{]}}}
\newcommand{\novpend}{{\color{red}\textbf{[NOVELTY PENDING]}}}

\newcommand{\lin}{\operatorname{lin}_e}
\newcommand{\Lead}{\operatorname{Lead}}
\newcommand{\val}{v}
\newcommand{\Lam}{\Lambda}
\newcommand{\qbin}[2]{\genfrac{[}{]}{0pt}{}{#1}{#2}_t}

\addbibresource{fpsac2027.bib}

\title[Valuations of Hikita's $\star$-product]{The $(s-1)$-adic valuation of Hikita's $\star$-product: a block law, cumulant leads, and a box symmetry}

\author{Rick (AI agent)\thanks{\todo{Robin: authorship, presenter and affiliation are your decision (asks sent 2026-10-05, no reply yet). The mathematics in this draft was produced by an AI agent; see the AI disclosure.}}\addressmark{1}}
\address{\addressmark{1}\todo{affiliation / contact}}

\received{\today}

\abstract{Hikita introduced a two-parameter deformation $\star$ of the product of symmetric functions. Writing $e_{\lambda_1}\star\cdots\star e_{\lambda_\ell}=\sum_\mu c_{\lambda\mu}(s,t)\,e_\mu$, we show that the $(s-1)$-adic valuation of $c_{\lambda\mu}$ is $\ell(\lambda)-\kappa(\lambda,\mu)$, where $\kappa$ is the maximal number of blocks into which $\lambda$ and $\mu$ split compatibly. Leading coefficients factor over block decompositions into connected pieces. For a full merge the connected piece is a connected-graph cumulant with edge weights $t^{\lambda_i\lambda_j}-1$. It interpolates between partition-lattice M\"obius values at $t=0$ and weighted Cayley counts at $t=1$, and a gauge-invariant ratio distinguishes it from Do{\l}{\k{e}}ga's column cumulants for $t\neq0$. The coefficients are invariant, up to a power of $s$, under complementation in a box. All second-order leading coefficients are given in closed form, via a residue formula for Green polynomials at two-part classes.}

\keywords{Hikita star product, Hall--Littlewood functions, elementary symmetric functions, cumulants, valuations}

\begin{document}

\maketitle

%% ===================================================================
\section{Introduction}
%% ===================================================================

Hikita~\cite{Hikita25} defined a commutative, associative product $\star$ on $\Lam_{q,t}$ by transporting the ordinary product through the level-one polynomial representation of the affine Hecke algebra; we always say \emph{Hikita's} $\star$. We write $s:=q^{-1}$. At $s=1$, $\star$ is the ordinary product, and at $s=0$ (after rescaling) it degenerates to a Hall--Littlewood product (\cref{rem:H}). Di Francesco and Kedem~\cite{DFK19} observed that the associated $t$-deformed graded characters lose Schur positivity at generic $t$, and positivity of $\star$-structure constants fails in every natural basis we tried. \emph{Positivity is lost, but structure is not.} This abstract shows that the structure survives in the $(s-1)$-adic filtration.

We expand the structure constants of $e^\star_\lambda:=e_{\lambda_1}\star\cdots\star e_{\lambda_\ell}=\sum_\mu c_{\lambda\mu}(s,t)\,e_\mu$ around $s=1$. The point of departure is Hikita's formula for $e_k\star(\cdot)$ as a sum over $k$-subsets (\cref{prop:subset}). In it, $s$ enters only through an Euler operator, so the $(s-1)^p$ Taylor piece of $e_k\star(\cdot)$ is a differential operator of order $p$, which can glue at most $p$ existing blocks of parts to a new part. This gives a lower bound on the valuation; exactness is settled at $t=0$, where $e^\star_\lambda$ is a modified Hall--Littlewood function and no cancellation occurs. Main results:
\begin{itemize}
  \item \textbf{Block valuation law} (\cref{thm:blocklaw}): $\val_{s-1}(c_{\lambda\mu})=\ell(\lambda)-\kappa(\lambda,\mu)$ for all $\mu\trianglerighteq\lambda$, $\mu\neq\lambda$, where $\kappa$ is the maximal number of compatible blocks. The first-order term is a symmetric biderivation with an explicit matrix (\cref{thm:bider,prop:M}).
  \item \textbf{Merge weights and full-merge leads} (\cref{thm:W,thm:G}): the leading coefficient for $\mu=(n)$ is $\frac{1-t^n}{\prod_i(1-t^{\lambda_i})}\sum_{H\text{ connected}}\prod_{ij\in H}(t^{\lambda_i\lambda_j}-1)$. The graph identity is due to Do{\l}{\k{e}}ga~\cite{Dolega19}; the identification with $\star$ is ours.
  \item \textbf{Block multiplicativity} (\cref{thm:F,thm:blockmult}): every leading coefficient is a sum over block decompositions of products of connected leading coefficients ($\exp$ of connected).
  \item \textbf{Box Complement} (\cref{thm:box}): $c_{N^\ell-\lambda,\,N^\ell-\mu}=s^{N\binom{\ell}{2}-(\ell-1)|\lambda|}\,c_{\lambda\mu}$, together with a Column Lemma.
  \item \textbf{Second order} (\cref{thm:v2,cor:v2}): every leading coefficient at valuation $2$ is in closed form, via a constant-term formula and two-part Green polynomials~\cite{Green55,DLT94,Macdonald95}.
\end{itemize}
At valuation $1$ the leads are governed by pairing with $p_n$, evaluated by a principal specialisation, i.e.\ one $t$-string $1,t,\dots,t^{k-1}$ (\cref{lem:linT}); at valuation $2$ by pairing with $p_xp_y$, where the poles of a constant-term integral fall into two interacting $t$-strings (\cref{thm:2pt}). The full-merge lead is a cumulant: up to a prefactor it is the connected part of the character $e_k\mapsto t^{\binom k2}$. This places $\star$ next to the Macdonald cumulants of Do{\l}{\k{e}}ga~\cite{Dolega19} and Haglund--Tewari~\cite{HaglundTewari26}. Their objects and ours share the graph half (\cref{thm:G} and the remark after it), but a gauge-invariant ratio separates $\star$ from Do{\l}{\k{e}}ga's $\oplus$-cumulants away from $t=0$.

\emph{Organisation.} \Cref{sec:bg} recalls Hikita's product and the dominance support. \Cref{sec:blocklaw} proves the block valuation law and the merge weights. \cref{sec:leads} treats leading coefficients, \cref{sec:box} the Box Complement theorem and the linear coefficient, and \cref{sec:v2} the second order. Proofs are sketched here; full proofs will appear in a longer version.

%% ===================================================================
\section{Background: Hikita's product and Hall--Littlewood notation}\label{sec:bg}
%% ===================================================================

\paragraph{Notation.} $\Lam=\mathbb{Q}(s,t)[e_1,e_2,\dots]$; $I=(e_1,e_2,\dots)$ is the augmentation ideal, so $I^b=\mathrm{span}\{e_\mu:\ell(\mu)\ge b\}$. For $G$ homogeneous of degree $n$, $\lin G:=[e_n]G$. $[m]=[m]_t=(1-t^m)/(1-t)$, $\varphi_m=\prod_{i=1}^m(1-t^i)$, $n(\lambda)=\sum_i(i-1)\lambda_i$, $\trianglerighteq$ is dominance. $P_\lambda(x;t)$, $Q_\lambda=b_\lambda P_\lambda$, $Q'_\lambda$ are the Hall--Littlewood functions of~\cite[Ch.~III]{Macdonald95}, with the Hall pairing $\langle\cdot,\cdot\rangle$ and the HL pairing $\langle p_\lambda,p_\mu\rangle_t=\delta_{\lambda\mu}z_\lambda\prod_i(1-t^{\lambda_i})^{-1}$.

\paragraph{Hikita's product.} $F\star G:=\mathfrak{q}(\mathfrak{q}^{-1}F\cdot\mathfrak{q}^{-1}G)$, where $\mathfrak q$ is the level-one polynomial representation map $F(Y)\mapsto F(Y)\bullet 1$ of the affine Hecke algebra of $GL_m$. By \cite[Lemma~3.1, Cor.~3.9]{Hikita25}, $\mathfrak q$ is a $\mathbb{Q}(q,t)$-linear isomorphism, compatible with $m\to m-1$, so $\star$ is well defined and restricts to $\Lam$ \cite[Def.~3.4, Lemma~3.2, Cor.~3.10]{Hikita25}. All structure constants below lie in $\mathbb{Q}[s,t]$ (\cref{thm:DS}); specialisations in $s$ are specialisations of these. We use $s=q^{-1}$; then Hikita's Pieri rule \cite[Thm.~3.12]{Hikita25} reads $e_1\star e_r=(1-s)[r+1]e_{r+1}+s\,e_1e_r$.

% registry: subset-formula-Ak-Kk, grade: proved
\begin{proposition}[Subset formula]\label{prop:subset}
For $F\in\Lam_N=\mathbb{Q}(s,t)[x_1,\dots,x_N]^{S_N}$, $N\ge\deg$,
\[ e_k\star F=E_kF:=\sum_{|A|=k}c_A\,X_A\,F(X_{A^c},sX_A),\qquad c_A=\prod_{i\in A,\,j\notin A}\frac{x_i-tx_j}{x_i-x_j},\quad X_A=\prod_{i\in A}x_i. \]
\end{proposition}
\begin{proof}[Proof idea]
By \cite[Def.~3.4, Lemma~3.3]{Hikita25}, $e_k\star F=t^{-\binom k2}e_k(Y)\bullet F$, with $Y_i$ the Cherednik operators in $N$ variables. Write $e_k(Y)$ with increasing indices, and let $\pi F=x_1F(x_2,\dots,x_N,sx_1)$. A Coxeter length count on minimal coset representatives of $S_k\times S_{N-k}$ gives $e_k(Y)F=t^{\binom k2}\sum_{|D|=k}T_{w^D}\,\pi^kF$. For $G$ symmetric in the first $k$ and the last $N-k$ variables, the chain factorisation of this coset sum and the Poincar\'e identity $\sum_{w\in S_k}w\bigl(\prod_{i<j}\frac{x_i-tx_j}{x_i-x_j}\bigr)=[k]_t!$ give $\sum_DT_{w^D}G=\sum_{|A|=k}c_A\,G^{(A)}$, where $G^{(A)}$ places the head variables on $A$.
\end{proof}

Write $e^\star_\lambda:=E_{\lambda_1}\cdots E_{\lambda_\ell}(1)=\sum_\mu c_{\lambda\mu}(s,t)\,e_\mu$. The variable $s$ enters only through $F(X_{A^c},sX_A)=s^{\Delta_A}F$, $\Delta_A=\sum_{i\in A}x_i\partial_i$; Taylor expansion at $s=1$ gives
\begin{equation}\label{eq:taylor}
E_k=\sum_{p\ge0}(s-1)^pE_k^{(p)},\qquad E_k^{(p)}=\sum_{|A|=k}c_AX_A\binom{\Delta_A}{p},
\end{equation}
and $E_k^{(p)}$ is a differential operator of order $\le p$ on $\Lam_N$, with $E_k^{(0)}=e_k\cdot$ and $E_k^{(p)}(1)=0$ for $p\ge1$.

% registry: full-upset-support-exact-valuation (root: DS), grade: proved
\begin{theorem}[Dominance support; the $s=0$ valuation]\label{thm:DS}
$e^\star_\lambda=s^{n(\lambda)}e_\lambda+\sum_{\mu\triangleright\lambda}c_{\lambda\mu}e_\mu$ with $c_{\lambda\mu}\in\mathbb{Q}[s,t]$ and $c_{\lambda\mu}|_{s=1}=0$. Moreover $c_{\lambda\mu}\neq0$ for every $\mu\trianglerighteq\lambda$, and the $s$-adic valuation of $c_{\lambda\mu}$ is exactly $n(\mu)$. At $s=t=0$, in the basis $b_\mu=s^{n(\mu)}e_\mu$, $e^\star_\lambda$ is the zeta function of dominance order.
\end{theorem}
\begin{proof}[Proof idea]
A termwise degree count on the subset formula: for symmetric $P$ let $D_j(P)$ be the maximal total degree in $j$ variables; $E_k$ raises $D_j$ by at most $\min(k,j)$, and $D_j(P)\le\sum_i\min(\rho_i,j)$ for all $j$ characterises $\mathrm{span}\{e_\nu:\nu\trianglerighteq\rho\}$ \cite[I (1.11), (6.6)]{Macdonald95}.
\end{proof}

% registry: theorem-H-s0-star-is-HL-pieri + d-equals-hall-littlewood-transition, grade: proved (Clio review 2026-10-06: H novel-as-checked, folklore-adjacent; d = known object)
\begin{remark}[The $s\to0$ edge; no novelty claim for the matrix]\label{rem:H}
Since $s=q^{-1}$, the edge $s\to0$ is Hikita's $q\to\infty$, where the unrescaled product degenerates \cite[Lemma~6.3]{Hikita25}. In the basis $b_\mu$, $e_k\star(\cdot)$ preserves $\bigoplus_\mu\mathbb{Q}[s,t]\,b_\mu$, and modulo $s$ it is conjugate, via $b_\mu\mapsto t^{-n(\mu')}P_{\mu'}(x;t^{-1})$, to multiplication by $t^{-\binom k2}e_k$. Consequently the leading coefficients $d_{\lambda\mu}(t):=[s^{n(\mu)}]\,c_{\lambda\mu}$ of \cref{thm:DS} \emph{are the known $e\to$HL transition matrix} of \cite[\S3.2]{Kirillov98}: writing $e_\lambda=\sum_\kappa M(e,P)_{\lambda\kappa}P_\kappa(x;t)$,
\[ d_{\lambda\mu}(t)=t^{\,n(\mu')-n(\lambda')}\,M(e,P)_{\lambda\mu'}(t^{-1}). \]
Their positivity and $d_{\lambda\mu}(1)=\#\{0\text{-}1\text{ matrices with row sums }\lambda,\text{ column sums }\mu'\}$ are therefore classical \cite[\S3.2]{Kirillov98}; what is new is only the identification of this matrix with the $s$-leading coefficients of $\star$.
\end{remark}

%% ===================================================================
\section{The block valuation law at $s=1$}\label{sec:blocklaw}
%% ===================================================================

\begin{definition}
For $\lambda,\mu\vdash n$, $\kappa(\lambda,\mu)$ is the largest $r$ such that $\lambda=\lambda^1\sqcup\cdots\sqcup\lambda^r$ and $\mu=\mu^1\sqcup\cdots\sqcup\mu^r$ as multisets of parts, with every block nonempty and $\mu^i\trianglerighteq\lambda^i$ (so $|\mu^i|=|\lambda^i|$); $\kappa=-\infty$ if $\mu\not\trianglerighteq\lambda$. $\mu$ is a \emph{coarsening} of $\lambda$ if $\kappa(\lambda,\mu)=\ell(\mu)$. $\val$ denotes the $(s-1)$-adic valuation.
\end{definition}

% registry: s1-first-order-star-is-biderivation, grade: proved
\begin{theorem}[First order is a biderivation]\label{thm:bider}
Let $D_k:=E_k^{(1)}$ and $B(f,g):=\partial_s(f\star g)|_{s=1}$. Then $D_k$ is a derivation of $\Lam_N$, and
\[ B(f,g)=\sum_{k,l}M_{kl}\,\frac{\partial f}{\partial e_k}\frac{\partial g}{\partial e_l},\qquad M_{kl}:=D_k(e_l)=B(e_k,e_l)=M_{lk}. \]
So $B$ is a symmetric biderivation, the carr\'e du champ of $L=\frac12\sum M_{kl}\partial^2/\partial e_k\partial e_l$.
\end{theorem}
\begin{proof}[Proof idea]
$\binom{\Delta_A}{1}=\Delta_A$ is a derivation and left multiplication by $c_AX_A$ preserves this. Differentiate $f\star g=\sum_\rho a_\rho(s)E_\rho g$ at $s=1$ and use $D_k(1)=0$. Symmetry of $M$ is commutativity of $\star$ \cite{Hikita25}.
\end{proof}

% registry: B-matrix-M_kr-closed-form, grade: proved
\begin{proposition}[The matrix $M$]\label{prop:M}
For $k\ge r\ge0$, with $L(a,b):=\dfrac{(1-t^{a+b})(t^{ab}-1)}{(1-t^a)(1-t^b)}$,
\[ M_{kr}=B(e_k,e_r)=r\,e_ke_r+\sum_{j=1}^{r}L(k-r+j,\,j)\,e_{k+j}e_{r-j}. \]
Consequently $[(s-1)^1]e^\star_\lambda=\sum_{i<j}M_{\lambda_i\lambda_j}\,e_{\lambda\setminus\{\lambda_i,\lambda_j\}}$.
\end{proposition}
\begin{proof}[Proof idea]
Differentiate \cref{cor:pieri} (whose proof does not use \cref{prop:M}) at $s=1$. Since $\pi_a(b)|_{s=1}=\delta_{b0}$, only $y=r$ survives at $s=1$, and its derivative is $r\,e_ke_r$. For $y=r-j$, $j\ge1$, put $a=k-r+j$. Then $\partial_s\pi_a(j)|_{s=1}=[u^j]\sum_{m<a}\frac{ut^m}{1+ut^m}=(-1)^{j-1}[a]_{t^j}$, so $\partial_sC_{a,j}|_{s=1}=-\frac{[a+j]}{[a]}[a]_{t^j}=L(a,j)$. The second statement follows from \cref{thm:bider}, since $D_k(1)=0$.
\end{proof}

% registry: s1-block-valuation-law-day220 (Thm A) + thmC-block-law-exact (Thm C), grade: proved
\begin{theorem}[Block valuation law]\label{thm:blocklaw}
For all $\lambda,\mu\vdash n$ and all $t$, $\val(c_{\lambda\mu})\ge\ell(\lambda)-\kappa(\lambda,\mu)$. For $\mu\trianglerighteq\lambda$, $\mu\neq\lambda$, equality holds over $\mathbb{Q}(t)$ and at $t=0$:
\[ \val(c_{\lambda\mu})=\ell(\lambda)-\kappa(\lambda,\mu). \]
At $t=0$ the leading coefficient is $(-1)^{\ell-\kappa}N(\lambda,\mu)$ with $N(\lambda,\mu)\ge1$ a count of minimal raising-operator configurations.
\end{theorem}
\begin{proof}[Proof idea]
\emph{Lower bound} (subset formula only): by polarization, an order-$p$ operator glues at most $p$ existing ``blocks'' to the new part, so the $(s-1)^m$ coefficient is a sum of products over $\ge\ell-m$ blocks, each of type $\lambda_C$ by the degree count of \cref{thm:DS}. \emph{Equality}: at $t=0$, $\omega e^\star_\lambda=\widetilde H_\lambda(x;s)=s^{n(\lambda)}Q'_\lambda(x;s^{-1})$, and Macdonald's raising operators $Q'_\lambda(x;u)=\prod_{i<j}\frac{1-R_{ij}}{1-uR_{ij}}h_\lambda$ \cite[III (2.15$'$) with \S5 Ex.~7(a)]{Macdonald95} give $R_{ij}^k$ the coefficient $(u-1)u^{k-1}=-(s-1)s^{-k}$ at $u=s^{-1}$. So every transfer edge costs exactly one $(s-1)$ with sign $-1$, so nothing cancels; the minimum edge count is $\ell-\kappa$ (components $\le\kappa$; a path tree on each indecomposable block). Specialise $t=0\to$ generic.
\end{proof}
\begin{remark}[Input at $t=0$; (N)-freeness]\label{rem:t0input}
The $t=0$ edge $e^\star_\mu|_{t=0}=s^{n(\mu)}P_{\mu'}(x;1/s,0)=\omega\widetilde H_\mu(x;s)$ is \textbf{not ours}: it is assembled from \cite{DFK18} eq.~(5.15) (with $M_{k,1}=E_k|_{t=0}$), Cor.~5.8 in its level-one form (5.25), Cor.~5.18 (eq.~(5.27)) (numbering of arXiv:1505.01657v2), and \cite[VI (5.1), (4.14)(iv)]{Macdonald95}. In particular \cref{thm:blocklaw} does not depend on Cherednik's nonsymmetric theory.
\end{remark}
\begin{remark}[Novelty]
We have not found the block law stated anywhere. Its $t=0$ case is classical in substance: the raising-operator formula follows from \cite[eq.~(11)]{DLT94} and Jacobi--Trudi, after which the valuation count above is a short exercise. What is new is the lower bound for all $t$ (from the subset formula alone), the exactness over $\mathbb{Q}(t)$, and the merge weights below.
\end{remark}

% registry: thmB-coarsening-exact-valuation, grade: proved
\begin{theorem}[Coarsenings: merge-history expansion]\label{thm:coarse}
Let $\mu\neq\lambda$ be a coarsening of $\lambda$ and $m=\ell(\lambda)-\ell(\mu)$. Then $\val(c_{\lambda\mu})=m$, and $\Lead_{\lambda\mu}(t):=[(s-1)^m]c_{\lambda\mu}$ equals
\[ \Lead_{\lambda\mu}(t)=\sum_{\text{tight histories }H\text{ ending at }\mu}\ \prod_{\text{merges in }H}W_k(J),\qquad \Lead_{\lambda\mu}(0)=(-1)^m\#\{H\}\neq0, \]
where a history processes $\lambda_\ell,\dots,\lambda_1$, each part $k$ either opening a block or merging $p\ge1$ existing blocks of sizes $J=(j_1,\dots,j_p)$, and $W_k(J):=\lin\,[\cdots[E_k^{(p)},e_{j_1}],\dots,e_{j_p}](1)$.
\end{theorem}
\begin{proof}[Proof idea]
Pass to $\mathrm{gr}_I\Lam$: non-tight terms land in $I^{\ell-m+1}$; tight steps are multilinear derivations, so only linear parts $g_i\equiv\gamma_ie_{n_i}\bmod I^2$ matter. At $t=0$, a roots-of-unity evaluation, Weyl orthogonality and Murnaghan--Nakayama give $W_k(J)|_{t=0}=(-1)^p$. This proof does not use (N).
\end{proof}

% registry: thmW-merge-weight-closed-form (second proof: thmW-second-proof-KF-free = (KF)-free second proof, shares Lemma 1.4 (per Clio); NOT independent), grade: proved
\begin{theorem}[Merge weights]\label{thm:W}
With $n=k+|J|$ and $p=|J|$,
\[ W_k(J)=(-1)^{p}\,\frac{[n]_t}{[k]_t}\prod_{j\in J}[k]_{t^{j}}. \]
Consequently $(-1)^m\Lead_{\lambda\mu}(t)>0$ for coarsenings and every real $t>0$.
\end{theorem}
\begin{proof}[Proof idea]
$W_k(J)=(-1)^{d-p}\lin T_k(p_J)$ with $T_kg:=\sum_{|A|=k}c_AX_Ag(X_A)$. The parabolic factorisation $T_kP_\rho=P_{\rho+1^k}$ \cite[III (2.2)]{Macdonald95} and $\lin P_\lambda=(-1)^{n-k}\frac{[n]}{[k]}P_{\lambda-1^k}(1,t,\dots,t^{k-1})$ (from \cite[III.7 Ex.~2, III.2 Ex.~1]{Macdonald95}) give the ``discrete HL measure = principal specialisation'' identity $\lin T_k f=(-1)^d\frac{[n]}{[k]}f(1,t,\dots,t^{k-1})$ (\cref{lem:linT}); apply it to $f=p_J$.
\end{proof}
% registry: thmW-second-proof-KF-free (Day 223 Thm 1.5), grade: proved
\begin{lemma}\label{lem:linT}
For $f\in\Lam_k$ homogeneous of degree $d$ and $n=k+d$: \[\lin T_k(f)=(-1)^d\frac{[n]_t}{[k]_t}f(1,t,\dots,t^{k-1}).\]
\end{lemma}
\begin{remark}[Novelty]
\Cref{lem:linT} is a short corollary of \cite[III.7 Ex.~2, III (2.2)]{Macdonald95} and is probably known.
\end{remark}

%% ===================================================================
\section{Leading coefficients: graph cumulants and blocks}\label{sec:leads}
%% ===================================================================

% registry: conjG-full-merge-lead-connected-graph, grade: proved
\begin{theorem}[Full-merge lead]\label{thm:G}
For $\lambda\vdash n$ with $\ell=\ell(\lambda)\ge2$,
\[ \Lead_{\lambda,(n)}(t)=\frac{1-t^n}{\prod_i(1-t^{\lambda_i})}\,K_\lambda(t),\qquad K_\lambda(t):=\sum_{H\text{ connected on }[\ell]}\ \prod_{ij\in H}\bigl(t^{\lambda_i\lambda_j}-1\bigr). \]
$K_\lambda$ is the connected part (joint cumulant) of $t^{e_2(\lambda)}=\prod_{i<j}(1+(t^{\lambda_i\lambda_j}-1))$.
\end{theorem}
\begin{proof}[Proof idea]
(1) Tight histories ending in one block are increasing trees on $[\ell]$ (parent = the merging part). (2) With $N_v$ the $\lambda$-mass of the subtree at $v$, \cref{thm:W} telescopes along the tree to $\frac{1-t^n}{\prod(1-t^{\lambda_i})}\prod_{v\to c}(t^{\lambda_vN_c}-1)$. (3) Tree--graph identity: $\sum_{\text{incr.\ trees}}\prod_{v\to c}(\prod_{j\in T_c}x_{vj}-1)=\sum_{\text{conn.\ }H}\prod_{e\in H}(x_e-1)$ for generic $x$, by the minimum-vertex/component decomposition; specialise $x_{ij}=t^{\lambda_i\lambda_j}$.
\end{proof}
% registry: conjG-full-merge-lead-connected-graph (corollaries; corollary_t1 field), grade: proved
\begin{corollary}\label{cor:G}
(a) $\lambda=1^n$: $\Lead=(-1)^{n-1}[n]_t\,I_n(t)$ with $I_n$ the tree inversion enumerator \cite{MallowsRiordan68}. (b) $t=0$: $\Lead=(-1)^{\ell-1}(\ell-1)!$. (c) $t\to1$: $\Lead\to(-1)^{\ell-1}n^{\ell-1}$ (weighted Cayley). (d) $\operatorname{sgn}\Lead=(-1)^{\ell-1}$ for all $t>0$, $t\neq1$.
\end{corollary}
\begin{proof}[Proof of (b), (c)]
At $t=0$ every edge weight is $-1$ and $\sum_{H\text{ conn.}}(-1)^{|H|}=(-1)^{\ell-1}(\ell-1)!$. As $t\to1$, with $\varepsilon=t-1$, only spanning trees contribute to the $\varepsilon^{\ell-1}$ term of $K_\lambda$, and $\sum_T\prod_{ij\in T}\lambda_i\lambda_j=(\prod_i\lambda_i)\,n^{\ell-2}$ (weighted Cayley); the prefactor is $(-n\varepsilon)/((-\varepsilon)^\ell\prod\lambda_i)\cdot(1+O(\varepsilon))$.
\end{proof}
\begin{remark}[Prior art --- graph half]
The tree, connected-graph and cumulant forms of $K_\lambda$ are a special case of Do{\l}{\k{e}}ga \cite[Prop.~2.1 and (11), Lemma~2.3]{Dolega19} (Tutte polynomial $T(1,t)$ of the graph $M_\lambda$), cf.\ also \cite{GesselSagan96,MallowsRiordan68}; step (3) is a Penrose-type tree--graph identity \cite{Penrose67}. \textbf{Must cite, not claim.} Ours: the identification of the $\star$-lead with prefactor$\cdot K_\lambda$ (\cref{thm:coarse,thm:W,thm:G}) and \cref{thm:F,thm:blockmult}. Haglund--Tewari's single-row cumulant \cite[Def.~7.1, Thm.~7.3]{HaglundTewari26} (their $q$ is our $t$) is the same graph half: since $\langle\widetilde H_{(m)},e_m\rangle=q^{\binom m2}$, one has $\langle\kappa(\lambda),e_n\rangle=t^{n(\lambda')}(t-1)^{1-\ell}K_\lambda(t)$.
\emph{Separator.} Full-merge leads $L(\lambda)$ of any such theory are only defined up to the gauge $L(\lambda)\mapsto L(\lambda)\,a^{\ell(\lambda)}g(|\lambda|)\prod_if(\lambda_i)$ (rescaling generators, degree-wise normalisation, rescaling the expansion parameter). The first gauge invariant appears at $n=4$: $J:=L(1^4)L(2,2)/L(2,1,1)^2$ (the exponents of $a$, $g(4)$, $f(1)$ and $f(2)$ all cancel). For $\star$, \cref{thm:G} gives $J=\frac{(t^2+1)(t^3+3t^2+6t+6)}{(t+1)(t^2+t+2)^2}$, and Haglund--Tewari's $\langle\kappa,e_n\rangle$ gives the same value. Do{\l}{\k{e}}ga's $\oplus$-cumulant of columns (lowest $(q-1)$-coefficient of $[g_n]\widetilde H_{\lambda'}$, $g_k=\widetilde H_{1^k}$) gives, by computation, $J\equiv 3/2=J_\star(0)$, whereas $J_\star(1)=1$, the Cayley value. So up to gauge Do{\l}{\k{e}}ga's $\oplus$ is $t$-independent and agrees with $\star$ only at $t=0$, where the leads are partition-lattice M\"obius values (\cref{cor:G}(b)), while Haglund--Tewari's is the same graph object without $\star$.
\end{remark}

% registry: conjF-coarsening-lead-factorizes, grade: proved
\begin{theorem}[Coarsening leads factor]\label{thm:F}
For $\mu$ a coarsening of $\lambda$,
\[ \Lead_{\lambda\mu}=\sum_{\pi}\ \prod_{B\in\pi}\Lead_{\lambda_B,(|\lambda_B|)}, \]
summed over set partitions $\pi$ of the positions $[\ell]$ whose sorted block sums equal $\mu$ (singletons contribute $1$). There is no interleaving factor.
\end{theorem}
\begin{proof}[Proof idea]
Increasing forests $=$ set partition $\times$ increasing tree per block, and the telescoped weight is local to each tree.
\end{proof}

% registry: block-multiplicativity-thm61, grade: proved
\begin{theorem}[Block multiplicativity]\label{thm:blockmult}
Let $\kappa=\kappa(\lambda,\mu)\ge1$, $\ell=\ell(\lambda)$, $m=\ell-\kappa$. Then
\[ [(s-1)^m]\,c_{\lambda\mu}=\sum_{\pi}\ \sum_{(\nu^C)_{C\in\pi}}\ \prod_{C\in\pi}[(s-1)^{|C|-1}]\,c_{\lambda_C,\nu^C}, \]
where $\pi$ runs over set partitions of $[\ell]$ into exactly $\kappa$ blocks and $(\nu^C)$ over tuples with $\bigsqcup_C\nu^C=\mu$ and $\nu^C\trianglerighteq\lambda_C$; every factor then has $\kappa(\lambda_C,\nu^C)=1$. With \cref{thm:blocklaw} this is $\Lead_{\lambda\mu}$.
\end{theorem}
\begin{proof}[Proof idea]
Expand by polarization into histories; non-tight histories produce $>\kappa$ compatible blocks and vanish by maximality of $\kappa$; tight histories factor block by block and interleave uniquely. Uses commutativity of $\star$ \cite{Hikita25} only to pass from compositions to partitions.
\end{proof}
\begin{remark}
\Cref{thm:F} is the case where each $\nu^C$ is a single part. The factorisation has the shape of a Mayer (cluster) expansion, i.e.\ moment--cumulant inversion of the graph identity in \cref{thm:G}; the $\star$-specific content is that the connected leads are those of \cref{thm:G}, with no interleaving correction.
\end{remark}

%% ===================================================================
\section{Box Complement and the linear coefficient}\label{sec:box}
%% ===================================================================

% registry: box-complement-theorem, grade: proved
\begin{theorem}[Box Complement]\label{thm:box}
Let $\lambda,\mu$ have at most $\ell$ parts (padded with zeros) and all parts $\le N$; put $N^\ell-\lambda:=(N-\lambda_\ell,\dots,N-\lambda_1)$. Then
\[ c_{N^\ell-\lambda,\,N^\ell-\mu}(s,t)=s^{N\binom{\ell}{2}-(\ell-1)|\lambda|}\,c_{\lambda\mu}(s,t). \]
\end{theorem}
\begin{proof}[Proof idea]
Stability of the subset formula for every $N$ (also $N<$ degree), and an inversion lemma in $N$ variables: for $F=e_N^mG(1/x)$ with $G$ of degree $g$, \[E_{N-k}F=s^{m(N-k)-g}e_N^{m+1}(E_kG)(1/x)\] (substitute $y=1/x$, $B=A^c$). Iterate, and compare coefficients using the identity $e_{N-r}(x)=e_N(x)\,e_r(1/x)$ and \cref{thm:DS} (support has $\le\ell$ parts).
\end{proof}
% registry: box-complement-theorem (Cor 2.2, Cor 2.3), grade: proved
\begin{corollary}[Column Lemma; lead invariance]\label{cor:column}
$c_{\lambda+1^\ell,\,\mu+1^\ell}=s^{\binom{\ell}{2}}c_{\lambda\mu}$. Hence $\val$ and $\Lead$ are invariant under $(\lambda,\mu)\mapsto(\lambda+1^\ell,\mu+1^\ell)$ and $(\lambda,\mu)\mapsto(N^\ell-\lambda,N^\ell-\mu)$; e.g.\ $(2,2,2)\to(3,3)$ and $(2,2,2)\to(4,1,1)$ have the lead $[3]_t(t+2)$ of $(1,1,1)\to(3)$.
\end{corollary}
\begin{remark}[Prior art]
The inversion lemma follows in a few lines from \cite[Rem.~3.3]{DFK19} ($S\Gamma_IS=\Gamma_I^{-1}$ under $x\mapsto x^{-1}$) plus subset complement; stability is \cite[Prop.~3.8, Cor.~3.10]{Hikita25}. We have not found \cref{thm:box} or \cref{cor:column} for $\star$ in the literature.
\end{remark}

% registry: plethystic-lin-e-and-pieri-reproof (Thm 3.1), grade: proved
\begin{theorem}[Plethystic linear coefficient]\label{thm:lin}
For $G$ homogeneous of degree $d$: $\lin(e_k\star G)=(-1)^d\frac{[k+d]_t}{[k]_t}\,G[(s-1)[k]_t]$, where $G\mapsto G[(s-1)[k]_t]$ is the ring map $p_r\mapsto(s^r-1)[k]_{t^r}$.
\end{theorem}
% registry: plethystic-lin-e-and-pieri-reproof (Cor 3.2), grade: proved
\begin{corollary}[Pieri rule, all $s$]\label{cor:pieri}
$e_k\star e_r=\sum_{y=0}^{\min(k,r)}s^yC_{k-y,r-y}\,e_{k+r-y}e_y$ with $C_{a,0}=C_{0,b}=1$, $C_{a,b}=(-1)^b\frac{[a+b]}{[a]}\pi_a(b)$, $\pi_a(j)=[u^j]\prod_{m=0}^{a-1}\frac{1+sut^m}{1+ut^m}$.
\end{corollary}
\begin{proof}[Proof idea]
Support by \cref{thm:DS}; strip $y$ columns by \cref{cor:column} (valid for compositions) to reduce to $\lin E_{k-y}(e_{r-y})$; evaluate by \cref{thm:lin}, which itself is \cref{lem:linT} applied to the generating function $E_k(\prod_iE(u_i))=\prod_iE(u_i)\,T_k(\prod_i\prod_{a\in A}\frac{1+su_ix_a}{1+u_ix_a})$.
\end{proof}
\begin{remark}
\Cref{cor:pieri} was first obtained by a different route, a $k$-fold kernel computation; Hikita's \cite[Thm.~3.12]{Hikita25} is the case $k=1$. Explicit rules for $e_k\star e_\lambda$ with $\lambda$ of two or $\ell$ columns will appear in the long version.
\end{remark}

%% ===================================================================
\section{Second order: every valuation-$2$ lead}\label{sec:v2}
%% ===================================================================
% v7.5 gate lifted wake 228: cold recheck PASSED (Day 226: Thm 1.1, Lemmas 2.1--2.2, Thm 2.5, Thm 4.2 kill test; Day 227: Prop 2.3); Green wording per Day 227 §1.4.

By \cref{thm:blockmult}, valuation-$2$ leads reduce to connected leads with $\ell(\lambda)=3$, $\kappa=1$; by \cref{cor:column} to two-part $\mu$. Put $T_ag:=\sum_{|A|=a}c_AX_Ag(X_A)$ for $g\in\Lam_a$, and
\[ \Gamma_a(e_b,e_c):=E_a^{(2)}(e_be_c)-e_bE_a^{(2)}(e_c)-e_cE_a^{(2)}(e_b)=\sum_{1\le r\le b,\ 1\le q\le c}(-1)^{r+q}e_{b-r}e_{c-q}\,T_a(p_rp_q). \]

% registry: ell3-kappa1-reduction-and-lambda-contains-1 (Prop 5.1), grade: proved
\begin{proposition}\label{prop:reduce}
If $(a,b,c)$ is any ordering of $\lambda$ and $\kappa(\lambda,\mu)=1$, then $[(s-1)^2]c_{\lambda\mu}=[e_\mu]\bigl(\Gamma_a(e_b,e_c)+D_aD_b(e_c)\bigr)$.
\end{proposition}

% registry: ct-adjoint-formula-thm11, grade: proved
\begin{theorem}[Constant-term adjoint formula]\label{thm:CT}
For $g\in\Lam_a$, $F\in\Lam$: $\varphi_a(t)\,\langle T_ag,F\rangle_t=\mathrm{CT}\bigl[Z\,g(z)\,F(z_1^{-1},\dots,z_a^{-1})\,K(z)\bigr]$, where $Z=z_1\cdots z_a$ and $K(z)=\prod_{1\le i<j\le a}\frac{z_j-z_i}{z_j-tz_i}$, expanded in $z_i/z_j$ ($i<j$).
\end{theorem}
\begin{proof}[Proof idea]
$T_ag$ restricted to $a$ variables is $Zg$ and lies in $\mathrm{span}\{P_\mu:\ell(\mu)=a\}$. Evaluate $\langle G,q_\alpha\rangle_t=[y^\alpha]G$ via Cauchy \cite[III (4.2), (4.4), (4.8)--(4.9)]{Macdonald95}, write $\langle G,Q_\lambda\rangle_t$ as a constant term by the raising-operator formula \cite[III (2.15$'$)]{Macdonald95}, and symmetrise $K=\Delta\kappa$ with $\sum_w w(\kappa)=v_a(t)$ \cite[III (1.4)]{Macdonald95}; $Q_\lambda(z^{-1})$ is evaluated by \cite[III (2.1), (2.11)--(2.12)]{Macdonald95}. This is Jing's vertex-operator calculus \cite{Jing91} with the vertex operators hidden.
\end{proof}

% registry: two-point-string-formula-thm25 (+ shuffle-identity-lemma24), grade: proved
\begin{theorem}[Two-point formula]\label{thm:2pt}
Let $g\in\Lam_a$ be homogeneous of degree $d$, $n=a+d$, $x,y\ge1$, $x+y=n$, and $\Phi_a(g;x,y):=\langle T_ag,p_xp_y\rangle$. For $1\le A\le a-1$ put $B=a-A$ and $G_A(w)=\sum_kG_{A,k}w^k:=g(1,t,\dots,t^{A-1},w,wt,\dots,wt^{B-1})$. Then
\begin{multline*}
\Phi_a(g;x,y)=(-1)^a(1-t^x)(1-t^y)\Biggl\{\sum_{A=1}^{a-1}t^{-AB}\Bigl[\frac{G_{A,y-B}}{(1-t^A)(1-t^B)}\\+\frac{1}{1-t^a}\sum_{m\ge1}(t^{-Am}-t^{Bm})\,G_{A,y-B-m}\Bigr]\\
-\frac{g(1,t,\dots,t^{a-1})}{1-t^a}\sum_{j=0}^{a-1}t^{-jy}\Biggr\}.
\end{multline*}
\end{theorem}
\begin{proof}[Proof idea]
Insert $S(u)=\sum_i1/(uz_i-1)$ into \cref{thm:CT} and take iterated residues: every variable lands on a $t$-string $1/u,t/u,\dots$ or $1/v,t/v,\dots$. A string of length $A$ with base $q$ has weight $(-1)^{A-1}\varphi_{A-1}q^A$, and two strings interact through the shuffle identity
\[ \mathrm{Sh}_{A,B}(w)=\qbin{A+B}{A}\,t^{-AB}\,\frac{(1-w)(1-t^{B-A}w)}{(1-t^{-A}w)(1-t^Bw)},\qquad w=u/v, \]
proved by a last-letter recursion. Take Taylor coefficients at $w=0$.
\end{proof}
\begin{remark}[Relation to Green polynomials]
Write $p_\mu=\sum_\lambda X^\lambda_\mu(t)P_\lambda(x;t)$. For $g=P_\rho$ and $\lambda=\rho+1^a$ one has $\Phi_a(P_\rho;x,y)=(1-t^x)(1-t^y)X^{\lambda}_{(x,y)}(t)/b_\lambda(t)$, so \cref{thm:2pt} evaluates Green polynomials at two-part classes; the two $t$-strings are the Frobenius orbits of a maximal torus of type $(x,y)$ \cite{Green55}. \Cref{thm:CT} is Jing's constant-term expression for the matrix element $\langle Q_\lambda,p_\mu\rangle$ \cite{Jing91}, cf.\ \cite[(2.17), (2.20)]{JingLiu21}. Jing and Liu \cite[Thm.~2.7]{JingLiu21} expand the same matrix element as a nested sum valid for all $\mu$, extracting one variable at a time; at two-part $\mu$ that sum does not close up, since its inner sums run over classes of arbitrary length. Their closed forms \cite[Thm.~2.10, (2.37)--(2.40)]{JingLiu21} restrict the Hall--Littlewood index $\lambda$, not the class $\mu$; the case $\ell(\lambda)=2$ goes back to Morris \cite{Morris77}.\footnote{We have not seen \cite{Morris77} first-hand; its scope is as described in \cite[p.~12, before Thm.~3.2]{JingLiu21}.} We evaluate by residues instead, and at $\ell(\mu)=2$ the poles organise into two $t$-strings. The formula is explicit, with no sum over classes, and linear in the two-string specialisation $G_A$ of $g$; for general $\rho$, $G_A$ is itself a finite (branching) sum, so \cref{thm:2pt} is not a product formula.
\end{remark}
\begin{example}[Two rows]
For $a=2$, only $A=B=1$ occurs, and $G_1(w)=P_\rho(1,w;t)=w^{\rho_2}\frac{(1-tw)-w^m(w-t)}{1-w}$ with $m=\rho_1-\rho_2\ge1$ ($G_1=w^{\rho_2}$ if $m=0$). Substituting into \cref{thm:2pt} gives, for $\lambda=(\lambda_1,\lambda_2)$ and $y\le x$,
\[ X^{(\lambda_1,\lambda_2)}_{(x,y)}(t)=\begin{cases}(t-1)\,t^{\lambda_2-1-y}(1+t^{y}), & y<\lambda_2,\\ m_{xy}-(1-t)\,t^{\lambda_2-1}, & y=\lambda_2,\\ (t-1)\,t^{\lambda_2-1}, & y>\lambda_2,\end{cases} \]
in agreement with a Kostka--Foulkes computation for $|\lambda|\le10$. It also follows directly from $X^\lambda_\rho=\sum_\nu\chi^\nu_\rho K_{\nu\lambda}(t)$ \cite[III (7.6$'$)]{Macdonald95}: for $\lambda=(n-k,k)$ one has $K_{(n-j,j),\lambda}=t^{k-j}$ (one tableau, and $K$ is monic of degree $n(\lambda)-n(\nu)$ \cite[III (6.5)]{Macdonald95}), and $\chi^{(n-j,j)}_\rho=\pi_j-\pi_{j-1}$ with $\pi_j$ the number of $j$-subsets fixed by $\rho$, so $X^\lambda_{(x,y)}=\pi_k+(t-1)\sum_{j<k}t^{k-1-j}\pi_j$. This is the two-part-class case of the range $\ell(\lambda)=2$ treated by Morris. On the diagonal class $(x,y)=\lambda$ it gives $t^{\lambda_2}-t^{\lambda_2-1}+m_{xy}$, which for $\lambda_1>\lambda_2$ with $\lambda_2\ge3$ odd equals $-1$ at $t=-1$ and $1$ at $t=0$. A root in $(-1,0)$ rules out a product of cyclotomic factors and powers of $t$, and Green polynomials at two-part classes admit no product formula even for two rows, as observed independently in~\cite[Thm.~D]{Clio26}.
\end{example}

% registry: gprime-v2-class4-open (Day 225 Thm 4.2; node id historical, class 4 CLOSED), grade: proved, cold recheck PASSED Day 226 (+ Prop 2.3 Day 227)
\begin{theorem}[Closed $\ell=3$, $\kappa=1$ leads]\label{thm:v2}
Let $\lambda=(a,b,c)$ be any ordering of a length-$3$ partition, $\mu=(x,y)$ with $\kappa(\lambda,\mu)=1$, $n=x+y$, $m_{xy}=2$ if $x=y$ and $1$ otherwise. Put \[\Lambda_a(r,q):=\lin T_a(p_rp_q)=(-1)^{r+q}\frac{[a+r+q]}{[a]}[a]_{t^r}[a]_{t^q}\] and $U_a(b,c;x,y):=\bigl((-1)^n\Phi_a(p_bp_c;x,y)-\Lambda_a(b,c)\bigr)/m_{xy}$. Then
\begin{multline*}
[(s-1)^2]\,c_{\lambda\mu}=(-1)^{b+c}U_a(b,c;x,y)+\sum_{\substack{1\le r<b\\ \{b-r,\,a+r+c\}=\{x,y\}}}(-1)^{r+c}\Lambda_a(r,c)\\
+\sum_{\substack{1\le q<c\\ \{c-q,\,a+b+q\}=\{x,y\}}}(-1)^{b+q}\Lambda_a(b,q)+[e_xe_y]\,D_a(M_{bc}),
\end{multline*}
with $D_a$ the derivation $D_a(e_j)=M_{aj}$ of \cref{prop:M}.
\end{theorem}
\begin{proof}[Proof idea]
\Cref{prop:reduce}; in the expansion of $\Gamma_a$ only terms with $r=b$ or $q=c$ survive $[e_xe_y]$; single-$e$ parts come from \cref{lem:linT}, and $\langle G,p_xp_y\rangle=(-1)^n(m_{xy}[e_xe_y]G+\lin G)$ converts \cref{thm:2pt} into $U_a$.
\end{proof}
% registry: gprime-v2-class4-open (Cor 4.3) + block-multiplicativity-thm61, grade: proved (Thm 4.2 inputs rechecked Day 226/227)
\begin{corollary}\label{cor:v2}
Every lead at valuation $2$ ($\kappa=\ell-2$, any $\ell$) has a closed formula: by \cref{thm:blockmult} it is a sum of products of pair leads (\cref{prop:M}) and $\ell=3$, $\kappa=1$ leads (\cref{thm:G}, \cref{cor:column}, \cref{thm:v2}).
\end{corollary}
\begin{example}
$(3,3,3)\to(7,2)$: $\Lead=2t^{13}+3t^{12}+3t^{11}+6t^{10}+6t^9+6t^8+9t^7+6t^6+3t^5+9t^4+6t^3+3t+4$ (agrees with direct computation). Leads are \emph{not} positive in general: $(4,4,2)\to(7,3)$ has $\Lead=(t+1)(t^2+1)(2t^8+t^7+t^6+3t^4+t^3+t^2-t+2)$.
\end{example}
\paragraph{Computer checks.} All statements in this section were checked against an independent implementation of Hikita's subset formula: \cref{thm:v2} on all $16$ class-$4$ pairs with $n\le12$ and on $27$ pairs in all $3$ orderings with $n\le10$; \cref{thm:2pt} in $229$ cases ($a\le5$); and the Green-polynomial dictionary symbolically for $|\lambda|\le8$ and numerically for $|\lambda|\le9$.

%% ===================================================================
\section{Open problems}
%% ===================================================================
\begin{enumerate}
  \item \textbf{Valuation $\ge3$.} Connected leads with $\ell\ge4$, $\kappa=1$. Hierarchy: $v=1\leftrightarrow p_n\leftrightarrow$ one $t$-string (\cref{lem:linT}); $v=2\leftrightarrow p_xp_y\leftrightarrow$ two strings (\cref{thm:2pt}). Does the three-string shuffle sum factor pairwise?
  \item \textbf{Zero set at generic $t$.} At $t=-2$ seven pairs with $n=6$ have $\val>\ell-\kappa$; for coarsenings this is cancellation between histories, e.g.\ $\Lead_{(1,1,1),(3)}=[3]_t(2+t)$. Describe the zero set of $[(s-1)^{\ell-\kappa}]c_{\lambda\mu}$ in $t$.
  \item \textbf{Counting.} A formula for $N(\lambda,\mu)$ in \cref{thm:blocklaw}; the observed $\Lead_{\lambda,(n-1,1)}(1)=2n^2-6n+3$ (computed only).
  \item \textbf{Positivity.} Leads are not $t$-positive (Example above), so no ``$t$-count of minimal flows'' can hold. Is there a signed or cancellation-free model?
  \item \textbf{Cumulant link.} Up to gauge, Do{\l}{\k{e}}ga's column cumulants \cite{Dolega19} are the $t=0$ specialisation of the $\star$-leads (at least for $n=4$, via $J$). Is there a Macdonald-cumulant construction whose leads are the $\star$-leads for all $t$, i.e.\ cumulants of $e_k\mapsto t^{\binom k2}$?
\end{enumerate}
% The s->0 edge remark is now \cref{rem:H} (after \cref{thm:DS}), per Clio's 2026-10-06 verdict.

\acknowledgements{\todo{Robin; Clio (peer agent) for reviews; per Robin.}}

\printbibliography

\clearpage
\section*{AI disclosure}
\todo{Robin: edit wording and fill in your role below. Mandatory, uncapped, not counted toward the 12 pages.}
% Also fill in the corresponding SoftConf submission-form field (AI-use declaration) with a short version of this text.

\textbf{The mathematics in this extended abstract was produced by an AI agent.} ``Rick'' is an autonomous AI research agent built on Anthropic's Claude models (Claude Opus~5.5 in the final sessions; earlier sessions ran on earlier Claude versions \todo{Robin: confirm the list}). It works in a sandbox with Python/SymPy, \LaTeX, and arXiv/Semantic Scholar search, and keeps a written log of every session. Rick chose the questions, ran the computer experiments, conjectured the statements, found and wrote the proofs, did the literature and novelty searches, and drafted this text.
\begin{itemize}
  \item \textbf{Statements and proofs.} Every result not attributed to a reference was conjectured and proved by the agent: the dominance support, the biderivation and the matrix $M$, the block valuation law, the merge-history expansion and merge weights, the full-merge, coarsening and block-multiplicative lead formulas, the Box Complement theorem and Column Lemma, the plethystic linear coefficient and Pieri rule, the constant-term and two-point formulas, and the valuation-$2$ leads. Each proof was written in full in the agent's notes; the merge weights and the second-order results (\cref{sec:v2}) were in addition re-derived from scratch in a later session. The proofs sketched here are condensed from those notes.
  \item \textbf{Computer verification.} All theorems were checked in small cases by the agent's own Python and SymPy code against a direct implementation of Hikita's subset formula. Sub-agents (further instances of the same model) were used for some computations and for drafting parts of the \LaTeX; their outputs were checked by the agent.
  \item \textbf{Review by a second AI agent.} A separate AI research agent (``Clio'', also Claude-based, run independently) reviewed some inputs, but not all: the subset-formula inputs; the $s\to0$ identification and the location of the matrix $d$ in the literature (\cref{rem:H}); the structure of the two proofs of \cref{thm:W} (they share one lemma, so they are not independent); an independent check of the dominance support for $n\le7$; and a numerical cross-check of the Green-polynomial dictionary after \cref{thm:2pt}. The block valuation law, \cref{thm:G,thm:F,thm:blockmult,thm:box}, \cref{thm:CT,thm:2pt,thm:v2} have not been reviewed by another agent.
  \item \textbf{Literature.} Prior-art searches were done by the agent and its sub-agents. Wherever a result or mechanism turned out to be known, it is cited and not claimed (e.g.\ \cref{rem:H}, the remarks after \cref{thm:G,thm:box}). Textbook equation numbers in \cite{Macdonald95} were checked against the book.
  \item \textbf{Human role.} \todo{Robin: describe your role (e.g.\ supervision, submission, any human verification). As of this draft no proof has been fully checked by a human.}
\end{itemize}

\end{document}
